% Replacement for the first two paragraphs of Section 1 (Introduction), paper.tex lines 31 and 33.
% Paragraph 1 replaces the old four-contribution paragraph. Paragraph 2 replaces the
% MMD/HSIC/dCov definition paragraph and now carries the contributions.
% Keeps the existing \cite keys gretton2012kernel, gretton2008kernel, szekely2007dcov.
% Needs \newcommand{\todo}[1]{\textbf{[TODO: #1]}}. Box 1 is numbered by hand in guidance.tex.

\section{Introduction}
Maximum mean discrepancy (MMD) \citep{gretton2012kernel} is often used to flag a change in distribution. The Hilbert--Schmidt independence criterion (HSIC) \citep{gretton2008kernel} flags dependence between paired variables. Distance covariance \citep{szekely2007dcov} asks the same question as HSIC through pairwise distances. All three tests were introduced for independent and identically distributed observations. A nominal $5\%$ test should reject a true null about five times in 100 repeated samples. On persistent data that promise can fail if the critical value ignores the dependence. This paper asks when the promise fails under long memory, and what an analyst should do instead. For independence testing our answer is a three-step rule. Remove the seasonal mean and the seasonal scale from each series. Then test with HSIC and circular shifts. A circular shift realigns one full series against the other by a wrapped offset, and the shifted statistics form the reference distribution. Finally, expect residual risk from slow drift that the adjustment cannot remove. Box~1 at the end of this section states the rule, the evidence behind it, and its limits.

The paper makes three contributions. First, it measures false-alarm rates under known null models and shows that the critical-value rule matters as much as the statistic. At $n=2000$, split-sample MMD with uncentred wild multipliers rejects $27.1\%$ of 1,000 stationary fractional Gaussian noise (fGn) series at $H=0.8$. HSIC with subsampling rejects $9.2\%$ of 2,000 independent fGn pairs at $H_X=H_Y=0.7$. A Hermite-rank calculation explains why the one-series split and the two-series product fail at different memory levels. \todo{PROP-B: one sentence on the half-mean spread proposition, if it lands in Section 2}. Second, a factorial experiment separates memory from heavy tails, nonlinear dynamics, and seasonality. A seasonal scale shared by two independent series is a separate failure, specific to kernel tests. Without memory it drives HSIC with wild multipliers to $73.3\%$, while Pearson correlation stays at $4.7\%$. \todo{PROP-A: one sentence on the random-phase seasonal-mean proposition}. The same experiments show that circular shifts hold size under both seasonal forms and keep their power at matched size. These results yield the decision rule. Third, we apply the rule to separated Nasdaq trade-duration windows. There the null cannot be verified, so the rejection fractions are descriptive, not false-alarm rates. \todo{STOCK-D2: one sentence on how the fractions move under the smooth diurnal factor}. The two headline rates come from different nulls and different rules. They illustrate the problem and do not rank MMD against HSIC.
